Hygroscopic particles that take up water vapour before mineral dust does could stop the dust from forming the droplets in which it freezes, and so suppress ice formation in mixed-phase clouds; how strongly depends on how much condensed water the dust needs before it can freeze. We tested this in six expansions in the fan-stirred AIDAd cloud chamber during the iSKYLAB campaign, in which a South African soil dust (SDSA01) and Arizona test dust (ATD03) were cooled from about $-12$ to $-30\,^\circ$C, alone and externally mixed with 4000--10\,500\cmc\ NaCl. We simulated every experiment with a bin-microphysics model that starts from the measured humidity, loses water vapour to the walls at a single fitted rate that predicts the cloud onset to within 30\,s, takes the ice-nucleation active-site density of each dust from its dust-alone experiment with a correction for fall-out, and has its wall-conditioned air fitted to the liquid phase only. With NaCl, the optical counter saw about ten times fewer cloud droplets in three of the four mixtures, while the hygrometer saw the same excess of total water as with the dust alone, and in the model none of the dust surface was in cloud droplets: the NaCl took the vapour. Three findings follow. (1) Dust froze before activating to a droplet. If it had to activate first, the model forms 0.000--0.049\cmc\ of ice in the mixtures, against 0.34--0.94\cmc\ observed, and cannot remove the water vapour that the hygrometer saw disappear at glaciation; this is ruled out for both dusts. (2) For ATD03 the competition suppressed the ice. Compared at the same temperature, before glaciation, the mixtures formed 0.44 and 0.69 times the ice of a model in which unactivated dust freezes as well as immersed dust (0.51 and 0.93 per unit dust surface from the observations alone): unactivated ATD03 froze about half as well as in droplets. (3) For SDSA01 the effect was the opposite. The mixtures formed five to six times more ice per unit dust surface than the dust alone and glaciated three degrees warmer, although the same powder was used throughout. Our leading hypothesis is that hydration changes the availability or effectiveness of the ice-active sites of the soil dust, which may involve its organic or mineral--organic components; contact freezing by the many small NaCl droplets cannot be excluded. In the model, the NaCl also keeps the dust airborne in small haze particles instead of in droplets that fall out, which increases the ice by factors of 1.2--2.1. Whether hygroscopic aerosol suppresses ice formation on dust therefore depends on how well the dust freezes before it activates, and this differs between dusts.